\documentclass[a4paper,10pt]{article}
\usepackage[francais]{babel}
\usepackage{url}

\selectlanguage{francais}
\input{letter-icube-uds.sty}

\selectlanguage{francais}
\begin{document}


\name{St\'ephane \textsc{Genaud}}
\fonction{Ma\^itre de conf\'erences (HDR)}
\telephoneillkirch{+33 (0)3688 54542}
\email{genaud@unistra.fr}

\begin{letter}[Recommendation of Marc Frincu for USC PostDoc position\\
on Cloud Computing and Big Data Informatics]%
{à}{Prof. Viktor Prasanna}
 

I first met Marc  Fr\^incu in 2007 when he was a student  on an internship in my
team at the  INRIA research center of Nancy,  France. I was at that  time on the
leave from my university on a researcher position. Although we have not directly
worked  together at  that time,  I  know from  my colleagues  that he  perfectly
achieved  his objectives  during  his  internship, which  led  my colleagues  to
propose him to apply for a \emph{CORDI} PhD funding at INRIA.

Marc preferred to  return to Romania, where he  went on with a PhD  on a subject
related to scheduling applications for Grids. This research theme is close to my
owns, and I  noticed during the last years the  interesting results published by
Marc in conferences in the domain.

During last summer, I learned that Marc was searching for a PostDoc position and
I immediately proposed  him to join us,  to work on scheduling  for clouds. This
theme perfectly matched  Marc's professional objectives. Marc joined  us on last
October,       to      participate       in      the       national      project
SONGS\footnote{http://infra-songs.gforge.inria.fr/}  I  am   involved  in.   The
objective  is to  enrich the  existing SimGrid  simulation toolkit  with an  API
dedicated to the  simulation of cloud-based systems. Simultaneously  to the work
on  simulation, we  carry  on theoretical  studies on  the  impact of  different
scheduling strategies  on different types  of workloads, that may  correspond to
certain  activity  types,  or  certain application  models  (e.g  submission  of
independent  jobs, or  workflows,  etc).  These  scheduling  algorithms will  be
candidate solutions  that could  benefit from  our cloud  simulator development.
Ideally, our  simulation results will  be compared to  those produced in  a real
software framework, such as the one under development in our team.


\paragraph{Technical Contributions}
Since he  arrived, Marc has  actively participated in  all these tasks. He is
undoubtly able to work autonomously and bring new ideas that he has thoroughly
tested before sharing them. He prefers to write them directly in a format that
will be efficient to make a paper of it. Marc is hence very good at producing
material such as text and experimental results that can be discussed by the 
rest of the team. He is also open-minded and takes into account all remarks
and suggestions that may be made about his work.

\paragraph{Current Contributions}

In details, his contributions during his 9 first months with us are as follows:
\begin{itemize}
\item I asked him to recap works from  the past year that were citing one of our
  paper, and which aspects were not addressed  yet. He came up very shortly with
  a very  good survey,  which we have  reworked and extended,  and that  we have
  submitted to a journal in february.
\item He has  focused on stuying cloud provisioning  plus scheduling techniqaues
  for  workflows.  He  has been  working  on this  on  his own,  with only  some
  feedback from  me. The  preliminary study  has been  publish in  the Cloudflow
  workshop at IPDPS this year.
\item We  are working  to extend  this work by  a 50\%  increment for  a journal
  special issue.
\item Marc is also involved  in the development of  the cloud component  of the
  SimGrid simulation toolkit. He has good programming skills, at least in Java
  and in shell scripts that he uses to produce graphic plots visualizing experiments
  results. He has been very quick to build up from existing Java code and spot 
  existing bugs.  
\end{itemize}

\paragraph{Leadership Potential}
This cloud  component development task is  distributed on three sites  and three
groups of  people (some  modeling hypervisors  at a  lower level,  some modeling
whole systems such as  EC2).  I would not say Marc has been  a leader in the job
of creating some emulation among the different groups.  Not mastering french was
probably an obstacle,  since some conversations ended up in  french.  However, i
tink Marc feels more comfortable when he can handle the whole thing. 

\paragraph{Proposal Compositions}
Marc has been very active at composing proposal regarding its academic work. When 
he arrived, he was still expecting some review from a PostDoc proposal he had
submitted in Romania. The month after, he was reusing the most relevant parts
of this document to propose me to participate in a bi-national scientific
cooperation program, to be submitted in last january. We ran out of time to
achieve the sublission, but this shows that Marc is always taking initiatives
in this field. 

As  a baseline, I would say have  been  favourably  impressed  so  far  by  the
capabilities of Marc to quickly identify and summarize relevant research issues,
and then to write it in a very  fluent style.  He is both a very autonomous, and
hard-working person.   I hope Marc's stay  (one year) at or  institution will be
long enough to achieve the numerous objective I would like him to undertake.
Finally, i  would summarize by saying that Marc's  work is discreet but efficient.

\end{letter}

\begin{flushright}
\includegraphics[width=.30\textwidth]{signgenaud.jpg}
\end{flushright}
%\newpage
%\bibliographystyle{plain}
%\bibliography{julien_blaise_ED}


\end{document}

\end{letter}
\begin{flushright}
\includegraphics[width=.20\textwidth]{signgenaud.jpg}
\end{flushright}
%\newpage
%\bibliographystyle{plain}
%\bibliography{julien_blaise_ED}


\end{document}
